% LaTeX version of REPORT.md (generated from REPORT.backup.md content).
% Compile: pdflatex REPORT.tex  (requires hyperref, booktabs)
\documentclass[11pt,a4paper]{article}

\usepackage[utf8]{inputenc}
\usepackage[T1]{fontenc}
\usepackage{lmodern}
\usepackage{microtype}
\usepackage{geometry}
\geometry{margin=1in}
\usepackage{hyperref}
\usepackage{booktabs}
\usepackage{longtable}
\usepackage{enumitem}
\usepackage{tabularx}

\hypersetup{
  colorlinks=true,
  linkcolor=blue,
  urlcolor=blue,
  citecolor=blue,
}

\title{ZTF / ALeRCE Stamp Benchmark for LLM Astronomy\\(Early-Stage Report)}
\author{}
\date{}

\begin{document}
\maketitle

\noindent\textbf{Status:} dataset construction and asset packaging are complete; \textbf{zero-shot prompts} (\S3) are fixed for initial runs; quantitative evaluation and metrics are still to be filled in. This document follows the \textit{shape} of technical reports such as \textit{Instructing LLMs to Negotiate using Reinforcement Learning with Verifiable Rewards} (problem $\rightarrow$ setup $\rightarrow$ metrics) and the \textit{AstroAlertBench} / \textit{LLM for Astronomy} framing (multimodal ZTF alert triage, ALeRCE stamp classifier, five-class taxonomy), while staying brief outside the dataset section.

\medskip
\noindent\textbf{Repository:} \url{https://github.com/Cruuusade/LLM_FOR_ASTRONOMY}

\section*{Abstract}

We are building a \textbf{balanced multimodal benchmark} for vision--language evaluation on \textbf{Zwicky Transient Facility (ZTF)} alert cutouts brokered by \textbf{ALeRCE}, using the operational \textbf{stamp classifier} taxonomy: supernova (SN), active galactic nucleus (AGN), variable star (VS), asteroid, and bogus. For each class we retain \textbf{500} objects selected by \textbf{descending ALeRCE stamp-classifier probability} (with a minimum probability floor and \textbf{replacement} from the ranked pool when stamp retrieval fails). Each example includes \textbf{tabular metadata} from the object query and \textbf{three FITS stamps} (science, template, difference), plus \textbf{PNG montages} for LLM/VLM consumption. This report records what has been implemented, the \textbf{standardized multimodal prompt} (Parts A--C + JSON) for zero-shot evaluation, and \textbf{reserved slots} for \textbf{metrics} aligned with multimodal astronomy benchmarks (e.g.\ AstroAlertBench-style grounding and classification).

\section{Introduction and Problem Formulation}

Operational ZTF pipelines issue large numbers of \textbf{first-detection} alerts that must be triaged into scientifically useful categories. ALeRCE's \textbf{stamp classifier} maps first-alert image triplets to five high-level classes, providing a \textbf{scalable weak label} suitable for benchmarking general models against a broker that astronomers already use (\href{https://science.alerce.online/}{ALeRCE}). Our goal is not to reproduce the neural stamp network, but to \textbf{standardize inputs} (images + metadata) and \textbf{labels} (stamp-classifier class + probability statistics) so that \textbf{LLMs / VLMs} can be compared on the same evidence used in broker-style triage---consistent with the problem setup described in \textit{AstroAlertBench: Evaluating Vision Language Models for Multimodal Astronomical Alert Triage} (ZTF cutouts, ALeRCE taxonomy, metadata serialization).

At this \textbf{early stage} we prioritize: (i)~\textbf{reproducible data collection}, (ii)~\textbf{class balance} (500 per class), (iii)~\textbf{high-confidence sampling} with documented floors and \textbf{replacement} when APIs return errors instead of FITS.

\section{Dataset Construction (Primary Section)}

\subsection{Source and Labels}

\begin{tabularx}{\textwidth}{@{}l X@{}}
\toprule
\textbf{Item} & \textbf{Description} \\
\midrule
\textbf{Survey} & ZTF (via ALeRCE API) \\
\textbf{Classifier} & ALeRCE \texttt{stamp\_classifier} (broker real-time stamp classifier setting) \\
\textbf{Classes} & \texttt{SN}, \texttt{AGN}, \texttt{VS}, \texttt{asteroid}, \texttt{bogus} (five-way classification) \\
\textbf{Selection} & Per class: query objects with \texttt{order\_by=probability}, \texttt{order\_mode=DESC}, minimum probability $\geq 0.8$ (Explorer-equivalent filter). Take the top candidates until \textbf{500 successful stamp downloads} per class; if an object's stamps cannot be retrieved (e.g.\ HTTP returns JSON instead of FITS), \textbf{skip} and use the \textbf{next} object in the ranked list. \textbf{Object IDs are unique} across the whole benchmark (no duplicate \texttt{oid} across classes). \\
\bottomrule
\end{tabularx}

\subsection{Modalities per Example}

\begin{enumerate}[leftmargin=*]
  \item \textbf{Metadata (tabular):} Fields returned by ALeRCE \texttt{query\_objects} for each \texttt{oid} (e.g.\ coordinates, detection counts, times, \texttt{probability}, \texttt{classifier}, etc.), consolidated in \texttt{data/manifest.csv}.
  \item \textbf{Images (FITS):} For each accepted detection, three cutouts under \texttt{stamps/<class>/<oid>/}: \texttt{science.fits}, \texttt{template.fits}, \texttt{difference.fits} (63$\times$63 ZTF-style stamps).
  \item \textbf{Images (LLM-facing PNG):} One \textbf{montage} per object under \texttt{stamps\_llm/<class>/<oid>/montage.png}: three panels labeled \textbf{Science}, \textbf{Template}, \textbf{Image} (the third panel is the \textbf{difference} cutout; label text follows the montage script). Stretched to 8-bit with a percentile-based normalization for display.
\end{enumerate}

\subsection{Summary Statistics (Current \texttt{data/summary.json})}

Selection used \textbf{500} objects per class. Reported \textbf{min / max} stamp-classifier \textbf{probability} among the \textbf{selected} set (per class):

\begin{center}
\begin{tabular}{l r r r}
\toprule
\textbf{Class} & \textbf{Count} & \textbf{Min prob.} & \textbf{Max prob.} \\
 & & \textbf{(selected)} & \textbf{(selected)} \\
\midrule
SN & 500 & 0.896107 & 0.953128 \\
AGN & 500 & 0.856078 & 0.898750 \\
VS & 500 & 0.942687 & 0.968923 \\
asteroid & 500 & 1.000000 & 1.000000 \\
bogus & 500 & 0.951082 & 0.977276 \\
\bottomrule
\end{tabular}
\end{center}

\noindent\textbf{Note:} The minimum probability \textbf{within each class's 500} is the \textbf{lowest} score among that high-confidence slice (useful for reporting dataset difficulty). The global selection floor remains \textbf{0.8} (\texttt{min\_probability\_floor}).

\subsection{Artifacts and Reproducibility}

\begin{tabularx}{\textwidth}{@{}l X@{}}
\toprule
\textbf{Path} & \textbf{Role} \\
\midrule
\texttt{download\_alerce\_benchmark.py} & Fetch ranked pools; fill 500/class with replacement; write \texttt{data/manifest.csv}, \texttt{data/summary.json}, \texttt{data/replacement\_skips.log} \\
\texttt{build\_stamps\_llm\_montages.py} & FITS $\rightarrow$ PNG montages in \texttt{stamps\_llm/} \\
\texttt{data/manifest.csv} & One row per benchmark example (oid, class, probabilities, candid, paths, \ldots) \\
\texttt{stamps/}, \texttt{stamps\_llm/} & Large binaries; \textbf{not} tracked in git; regenerate locally \\
\bottomrule
\end{tabularx}

\subsection{What We Still Need (Dataset)}

\begin{itemize}[leftmargin=*]
  \item \textbf{Prompt-facing metadata subset:} Full broker metadata exists in principle; we still need a \textbf{fixed, documented} text serialization (which columns go into the model prompt vs.\ held out for scoring only)---as in AstroAlertBench's distinction between \textit{stored} vs \textit{prompt-facing} metadata.
  \item \textbf{Optional upgrade:} Align even more tightly with AstroAlertBench by attaching \textbf{AVRO-level} candidate fields for each chosen \texttt{candid} if we want parity with ``first-detection packet'' wording (current pipeline centers on object query + stamps + chosen candid).
  \item \textbf{Hosting:} Large FITS/PNG trees require \textbf{separate release} (e.g.\ Zenodo, Hugging Face datasets, or Git LFS) if public redistribution is desired; the GitHub repo is \textbf{code + manifests} only.
\end{itemize}

\section{Prompt Design (Zero-Shot)}

\noindent\textbf{Goal.} We construct a standardized multimodal prompt that combines image triplets (Science, Template, Difference) with selected metadata fields, and requires the model to produce a \textbf{structured, machine-parseable response} divided into \textbf{Part A (metadata grounding)}, \textbf{Part B (scientific interpretation)}, and \textbf{Part C (staged classification)}. This follows the evaluation protocol used in multimodal astronomy benchmarks such as AstroAlertBench.

\subsection{Prompt Structure}

\begin{itemize}[leftmargin=*]
  \item \textbf{Image inputs:} Science, Template, and Difference cutouts
  \item \textbf{Metadata:} Selected prompt-facing fields (e.g.\ band, magnitude, sgscore)
  \item \textbf{Instructions:} Task definition, class definitions, and output schema
\end{itemize}

\subsection{System Prompt}

\begin{quote}\small
You are an experienced astrophysicist. Your task is to classify astronomical transient candidates using three image cutouts (Science, Template, Difference) and associated metadata.

You must analyze the images and metadata carefully and produce a structured response divided into three parts:

\textbf{Part A: Metadata Grounding} \\
Extract and restate the key metadata fields and relevant observable properties.

\textbf{Part B: Scientific Interpretation} \\
Provide a concise scientific interpretation of the candidate based on the images and metadata.

\textbf{Part C: Staged Classification} \\
Make a step-by-step classification decision:
\begin{itemize}[leftmargin=*]
  \item Stage 1: Determine whether the detection is a real astrophysical object or bogus
  \item Stage 2: Determine whether it is astrophysical or non-astrophysical
  \item Stage 3: Assign one of the five classes: SN, AGN, VS, asteroid, or bogus
\end{itemize}

Your final answer must follow the required JSON format exactly.
\end{quote}

\subsection{User Prompt Template}

\noindent Fill placeholders \texttt{\{oid\}}, \texttt{\{fid / band\}}, \texttt{\{magpsf\}}, \texttt{\{sgscore1\}} (and any additional fields) from the manifest / AVRO for each alert.

\begin{quote}\small
\textbf{Alert ID:} \texttt{\{oid\}}

\textbf{Images:}
\begin{itemize}[leftmargin=*]
  \item Science image: recent observation
  \item Template image: reference observation
  \item Difference image: Science $-$ Template
\end{itemize}

\textbf{Metadata:}
\begin{itemize}[leftmargin=*]
  \item Band: \texttt{\{fid / band\}}
  \item Magnitude (magpsf): \texttt{\{magpsf\}}
  \item sgscore1: \texttt{\{sgscore1\}}
  \item Detection statistics and additional fields as provided
\end{itemize}

\textbf{Task:} Analyze the image triplet and metadata, then complete Parts A--C.

\textbf{Class Definitions:}
\begin{itemize}[leftmargin=*]
  \item \textbf{SN (Supernova):} Transient event not present in template, often point-like and located near a host galaxy
  \item \textbf{AGN:} Persistent or variably bright source associated with a galaxy nucleus
  \item \textbf{VS (Variable Star):} Stellar variability, typically present in both Science and Template images
  \item \textbf{Asteroid:} Moving object, may show positional shift or streak-like morphology
  \item \textbf{Bogus:} Artifacts such as noise, cosmic rays, subtraction errors, or misalignment
\end{itemize}

\textbf{Instructions:}
\begin{enumerate}[leftmargin=*]
  \item Focus on the central object in the images
  \item Compare Science, Template, and Difference images
  \item Use metadata to support your reasoning
  \item Follow the staged classification process
\end{enumerate}

\textbf{Output Format (strict JSON):}
\end{quote}

\begin{verbatim}
{
  "Part A": {
    "band": "...",
    "magnitude": "...",
    "sgscore": "...",
    "key_features": "..."
  },
  "Part B": {
    "interpretation": "..."
  },
  "Part C": {
    "stage1": "real_object | bogus",
    "stage2": "astrophysical | non_astrophysical",
    "stage3": "SN | AGN | VS | asteroid | bogus",
    "confidence": 0.0
  }
}
\end{verbatim}

\subsection{Design Rationale}

This structured prompt design enables:
\begin{itemize}[leftmargin=*]
  \item \textbf{Grounding evaluation:} Part A measures whether the model correctly uses metadata
  \item \textbf{Interpretability:} Part B separates reasoning from final prediction
  \item \textbf{Fine-grained evaluation:} Part C enables staged error analysis
  \item \textbf{Robust parsing:} JSON output ensures machine-readable results
\end{itemize}

\subsection{Future Extensions}

We plan to explore:
\begin{itemize}[leftmargin=*]
  \item Few-shot prompting with representative examples
  \item Ablations with and without metadata
  \item Alternative reasoning formats (e.g.\ chain-of-thought vs concise rationale)
\end{itemize}

\noindent Zero-shot and follow-up experiments will be run via the \textbf{Tinker} API.

\subsection{Reference: MeerLICHT Real/Bogus (Prior Work --- Binary Task)}

The following is \textbf{not} our ZTF five-class prompt; it is a reference template from the MeerLICHT / \href{https://github.com/turanbulmus/spacehack}{SpaceHack} line (binary Real vs Bogus). Our benchmark uses ZTF/ALeRCE naming and five classes. See SpaceHack materials for full persona and instructions.

\section{Evaluation Metrics \textit{(Reserved --- to be finalized)}}

Following the \textit{separation of concerns} in the negotiation RL paper (clear \textbf{metric definitions} before results) and AstroAlertBench's \textbf{Part A / B / C} style, we reserve the following \textbf{slots}. Actual numbers will be filled after running models.

\subsection{Classification Quality}

\begin{tabularx}{\textwidth}{@{}l X@{}}
\toprule
\textbf{Metric} & \textbf{Definition / notes} \\
\midrule
\textbf{Accuracy} & Fraction of correct top-1 class predictions vs.\ ALeRCE stamp label \\
\textbf{Macro-F1} & Unweighted mean of per-class F1 (important for balanced 500/ class) \\
\textbf{Per-class precision/recall} & Confusion matrix SN, AGN, VS, asteroid, bogus \\
\textbf{Calibration} & Reliability vs.\ model-reported confidence (if JSON includes probabilities) \\
\bottomrule
\end{tabularx}

\subsection{Format \& Grounding \textit{(optional, AstroAlertBench-inspired)}}

\begin{tabularx}{\textwidth}{@{}l X@{}}
\toprule
\textbf{Metric} & \textbf{Definition / notes} \\
\midrule
\textbf{JSON validity rate} & Parsable structured output rate (cf.\ ``Valid Rate'' in negotiation benchmarks) \\
\textbf{Metadata grounding} & When metadata is provided, score whether rationale cites correct fields \\
\textbf{Human / judge model review} & Subset scored for scientific plausibility of rationale \\
\bottomrule
\end{tabularx}

\subsection{Efficiency \& Cost}

\begin{tabularx}{\textwidth}{@{}l X@{}}
\toprule
\textbf{Metric} & \textbf{Definition / notes} \\
\midrule
\textbf{Tokens / USD per example} & For API-hosted VLMs \\
\textbf{Latency} & Wall time per batch \\
\bottomrule
\end{tabularx}

\section{Current Limitations}

\begin{itemize}[leftmargin=*]
  \item Labels are \textbf{broker stamp-classifier} outputs, not independent human truth; errors in ALeRCE propagate as label noise.
  \item \textbf{Asteroid} class in the current summary shows \textbf{probability 1.0} across the selected range---reporting and modeling should treat that \textbf{ceiling effect} explicitly.
  \item PNG montages use a \textbf{display stretch}; models should be evaluated with a \textbf{fixed preprocessing} policy declared in the final benchmark card.
\end{itemize}

\section{References \textit{(selected)}}

\begin{itemize}[leftmargin=*,label={}]
  \item ALeRCE: \url{https://science.alerce.online/}
  \item ZTF public data releases: \url{https://www.ztf.caltech.edu/ztf-public-releases.html}
  \item Carrasco-Davis et al.; ALeRCE stamp classifier (see ALeRCE publications).
  \item \textit{AstroAlertBench} manuscript (Claire Chen, Caltech) --- multimodal ZTF/ALeRCE alert triage benchmark framing.
  \item \textit{Instructing LLMs to Negotiate using Reinforcement Learning with Verifiable Rewards} (Claire Chen) --- example of abstract + numbered sections + explicit metric tables.
  \item MeerLICHT / SpaceHack prompting lineage: \url{https://github.com/turanbulmus/spacehack}
\end{itemize}

\medskip
\noindent\textit{Document version: early stage --- dataset and zero-shot prompt (\S3) are authoritative; quantitative metrics (\S4) are placeholders until experiments complete.}

\end{document}
